%% kapa: exam=1 round=31 year=2020 period=2 form=A subject=accounting subject_ko=회계학 seq=34 number=74
%% source: Q-Net 감정평가사 기출문제 (cst00309) artlSeq=5207219 file=extracted/1008/5207219/A형/A형/2020년도 제31회 감정평가사 1차 시험_2교시_A형.pdf
%% conversion: kordoc 4.24.51 -> tools/gen1.py (auto); review=proofread-sonnet; flags=none
\begin{question}{74}
(주)감평은 동일한 원재료를 결합공정에 투입하여 세 종류의 결합제품 A, B, C를 생산ㆍ판매하고 있다. 결합제품 A, B, C는 분리점에서 판매될 수 있으며, 추가가공을 거친 후 판매될 수도 있다. (주)감평의 20x1년 결합제품에 관한 자료는 다음과 같다.
\par{\small\centering\begin{tabular}{|c|c|c|c|c|}\hline
제품 & 생산량 & \makecell{분리점에서의\\단위당 판매가격} & 추가가공원가 & \makecell{추가가공 후\\단위당 판매가격} \\ \hline
A & 400단위 & ￦120 & ￦150,000 & ￦450 \\
B & 450단위 & 150 & 80,000 & 380 \\
C & 250단위 & 380 & 70,000 & 640 \\ \hline
\end{tabular}\par}
결합제품 A, B, C의 추가가공 여부에 관한 설명으로 옳은 것을 모두 고른 것은? (단, 기초 및 기말 재고자산은 없으며, 생산된 제품은 모두 판매된다.)
\begin{framedtext}
ㄱ. 결합제품 A, B, C를 추가가공 하는 경우, 단위당 판매가격이 높아지기 때문에 모든 제품을 추가가공 해야 한다. \par
ㄴ. 제품 A는 추가가공을 하는 경우, 증분수익은 ￦132,000이고 증분비용은 ￦150,000이므로 분리점에서 즉시 판매하는 것이 유리하다. \par
ㄷ. 제품 B는 추가가공을 하는 경우, 증분이익이 ￦23,500이므로 추가가공을 거친 후에 판매해야 한다. \par
ㄹ. 제품 C는 추가가공을 하는 경우, 증분수익 ￦65,000이 발생하므로 추가가공을 해야 한다. \par
ㅁ. 결합제품에 대한 추가가공 여부를 판단하는 경우, 분리점까지 발생한 결합원가를 반드시 고려해야 한다. \par
\end{framedtext}
\begin{choices}
  \item ㄱ, ㄴ
  \item ㄴ, ㄷ
  \item ㄱ, ㄴ, ㄷ
  \item ㄴ, ㄷ, ㄹ
  \item ㄷ, ㄹ, ㅁ
\end{choices}
\end{question}
\answer{2}
